\chapter{Conclusions}
\label{chap:conclusions}

Local robotic polishing requires a strategy for movement, contact loading, and time within a selected region. Human region specification supplies the spatial task, but does not by itself determine that temporal strategy. This thesis has addressed the resulting skill-generation problem through paired motion--force teaching, observation-conditioned sequence learning, and a common robot execution interface.

Tool position is learned relative to a frozen stain center and paired with orientation and target force on a shared sequence. The direction-transfer adapter aligns live XY and the policy image with the teaching direction, then restores generated XY by inverse rotation and center addition. Z, tool orientation, and force keep their conventions. The evaluated Flow Matching policy uses RGB, Cartesian state, and three-axis force history, with explicit coordinate restoration, pose conditioning, plan connection, and existing contact regulation. E1/E2 disable direction alignment and do not measure its transfer performance. The implementation separates upper-level skill generation, lower-level contact control, and completion recognition.

The first research question compares strategy providers for the fixed human-specified task. In the primary E1 cohort, rules, force-inclusive replay, and the conditioned policy each complete the required contact--loaded-pass--exit procedure in all five planned executions. Replay has a shorter mean completion time and smaller time dispersion than C, while C has greater mean excess-force exposure. The learned skill is therefore an executable strategy through the common robot interface, and rules and replay are also effective alternatives within this nominal task. C5 remains in the primary comparison; the shorter, lower-exposure C6 is a separate additional diagnostic.

The second research question compares the force-information configurations while retaining lower-level feedback. The force-observing motion--force policy C has a \BCSeconds~s (\BCPercent\%) shorter mean operator-based completion interval than force-masked B, with five completed attempts in each. Motion-only A completes four of five attempts; A4 does not progress to release and is interrupted. It remains in all-attempt force, travel, and elapsed-interval statistics, while A's successful-time mean uses four completions. Timing, phase progression, and force exposure differ across the configurations, without a monotonic improvement from A through B to C or an identified cause for the single interrupted transition.

The demonstrated scope is the generation and execution of a human-guided local contact skill on the marked-surface setup. It does not include autonomous recognition of defect severity, learned recognition of complete scratch removal, calibrated material-removal estimation, or measured surface restoration. The existing dose, blending, and adaptive-execution assets provide concrete extension paths, with their implementation and validation status stated separately from the main findings.

The thesis's contribution is the coherent connection from local task specification to a time-related motion--force representation, conditioned generation, shared execution, and corresponding evaluation. The fixed-task comparison establishes actual execution alongside valid rule and replay strategies; the force-information comparison establishes observed differences in completion, operator-based timing, phase progression, and force use. These findings answer the stated questions within the nominal contact-task scope.
